% ============================================================
% Method — README §3, §6, §7, §8
%   3.1 Problem formalization + carrier taxonomy (§3.2)
%   3.2 LLM-call boundary interception (§6.1)
%   3.3 Self-reported candidate sampling (§6.2 / §6.3)
%   3.4 Cryptographic audit trace (§8)
% ============================================================
\section{Method}
\label{sec:method}

% ---------- pipeline figure ----------
% ============================================================
% Figure: End-to-end MemMark pipeline (data flow diagram)
% ============================================================
\begin{figure*}[t]
\centering
\begin{tikzpicture}[
  node distance=6mm and 10mm,
  block/.style={draw, thick, rounded corners=2pt, minimum height=8mm, minimum width=22mm, align=center, font=\small, inner sep=3pt},
  sdk/.style={block, fill=blue!8},
  wrapper/.style={block, fill=orange!12},
  sampler/.style={block, fill=green!10},
  store/.style={block, fill=gray!10},
  audit/.style={block, fill=purple!8},
  flow/.style={-{Latex[length=1.4mm,width=1.2mm]}, thick},
  flowdash/.style={-{Latex[length=1.4mm,width=1.2mm]}, thick, dashed},
]

% Top row: ingestion path
\node[sdk] (event) {Conversation\\turn $e_t$};
\node[sdk, right=of event] (backend) {Memory SDK\\(\amem{} / \graphiti{})};
\node[wrapper, right=of backend] (wrap) {LLM-call\\wrapper};
\node[wrapper, right=of wrap] (enum) {Self-reported\\enumeration};
\node[sampler, right=of enum] (samp) {Binning\\sampler $\nonceT$};
\node[audit, below=10mm of samp] (commit) {Commitment\\$\commit$};
\node[audit, left=of commit] (merkle) {Per-session\\Merkle log};
\node[audit, left=of merkle] (anchor) {Signed anchor\\$\header_{T}$};

% Storage
\node[store, below=10mm of backend] (store) {Memory snapshot\\+ in-record sidecar};

% Top-row flow
\draw[flow] (event) -- (backend);
\draw[flow] (backend) -- node[above,font=\scriptsize]{prompt} (wrap);
\draw[flow] (wrap) -- node[above,font=\scriptsize]{+ instr.} (enum);
\draw[flow] (enum) -- node[above,font=\scriptsize]{$(\candset, \probdist)$} (samp);
\draw[flow] (samp.west) to[bend left=25] node[above,font=\scriptsize]{$c^*$} (wrap.east);

% Audit flow
\draw[flow] (samp) -- (commit);
\draw[flow] (commit) -- (merkle);
\draw[flow] (merkle) -- (anchor);

% Persistence
\draw[flowdash] (backend) -- (store);
\draw[flowdash] (commit) to[bend right=15] node[right,font=\scriptsize]{$\reveal_{t}$} (store);
\draw[flowdash] (anchor) to[bend left=15] node[below,font=\scriptsize]{$\header_T$} (store);

\end{tikzpicture}
\caption{\textbf{End-to-end \memmark{} pipeline.} The memory SDK
issues an internal LLM call (blue). Our wrapper (orange) appends an
\agentmark-style enumeration instruction; the LLM self-reports
$K$ alternatives with weights $(\candset, \probdist)$. The binning
sampler (green) keyed-picks $c^*$ via $\nonceT = \text{PRF}(K,
\ctx)$. Each pick produces a commitment, gets appended to the
per-session Merkle log, and contributes to the signed anchor
$\header_T$ (purple). At seal time the anchor and the reveal record
land directly inside the memory snapshot, enabling R3 in-record
verification without any external audit store.}
\label{fig:pipeline}
\end{figure*}


% ============================================================
\subsection{Problem Formalization and Carrier Taxonomy}
\label{sec:method-formal}
% ============================================================
\todo{Define memory state $M_t$, incoming event $e_t$, evolve
decision $\decision = \langle \carrier, \candset, \probdist, \ctx
\rangle$ with $\carrier \in \{\text{update}, \text{link},
\text{semantic}\}$. State the backend-invariance: \memmark{} reads
only $(\candset, \probdist, \ctx)$, never touches backend internals.
Match \S3.2 of the README. Include a table summarizing the three
carriers across \amem{} / \graphiti{}.}

% ============================================================
% Table: 3 carriers across A-Mem / Graphiti
% Lives in §3.1 Method.
% ============================================================
\begin{table}[t]
\centering
\small
\setlength{\tabcolsep}{4pt}
\begin{tabular}{lll}
\toprule
Carrier $\carrier$ & \amem{} instance & \graphiti{} instance \\
\midrule
\texttt{update\_target}
  & note id (which to update)
  & fact-edge id / episode \\
\texttt{link\_target}
  & keyword cluster / tag
  & entity attach point \\
\texttt{semantic\_realization}
  & note description + tags
  & edge label / fact phrasing \\
\bottomrule
\end{tabular}
\caption{\textbf{Three carriers, two backends}. Each $\carrier$ is a
non-trivial-candidate-set evolve decision realized in
backend-specific form; the watermark sampler reads only
$(\candset, \probdist, \ctx)$ and is therefore backend-invariant.}
\label{tab:carrier-taxonomy}
\end{table}


\todo{Then list the minimum adapter signature
($\texttt{enumerate\_candidates}$, $\texttt{score\_candidates}$,
$\texttt{apply\_selected}$) and the per-carrier reportable quantities
($|\candset|$, $H(\probdist)$, acceptance rate).}

% ============================================================
\subsection{LLM-Call Boundary Interception}
\label{sec:method-interception}
% ============================================================
\todo{Key observation: each SDK-internal LLM call is itself an
$\varepsilon$-equivalent decision point at $T{=}0.7$. We hot-swap
the SDK's LLM client with a wrapper (\amem{}: replace
\texttt{AgenticMemorySystem.llm\_controller}; \graphiti{}: replace
\texttt{Graphiti.llm\_client}). Wrapper appends an \agentmark-style
instruction to the SDK's prompt asking the LLM to enumerate $K$
alternatives with self-reported weights in a single call. The wrapper
parses, runs keyed pick via the binning sampler, returns the chosen
decision in the SDK's expected schema. Backend doesn't know.}

% prompt template figure
% ============================================================
% Figure: AgentMark-style prompt instruction (verbatim)
% ============================================================
\begin{figure}[t]
\centering
\begin{tcbox}[colback=gray!4, colframe=black!50, fontupper=\ttfamily\scriptsize, width=0.97\linewidth, boxsep=2pt, left=2pt, right=2pt, top=2pt, bottom=2pt]
CRITICAL OVERRIDE: instead of returning a single answer, return JSON in EXACTLY this multi-candidate form:\\[2pt]
\{\\
\hspace*{1em}"candidates": [\\
\hspace*{2em}\{"decision": <answer matching the original schema>, "weight": <float>\},\\
\hspace*{2em}\{"decision": <plausible alternative>, "weight": <float>\},\\
\hspace*{2em}...\\
\hspace*{1em}],\\
\hspace*{1em}"thought": "..."\\
\}\\[2pt]
- Provide K candidate alternatives;\\
- Each "decision" must independently match the original schema;\\
- Weights $>$ 0, rank by your real preference; sum normalized.
\end{tcbox}
\caption{\textbf{Open-vocabulary $\to$ discrete enumeration prompt.}
Appended verbatim to the memory SDK's original prompt.  Reduces an
unbounded structured-JSON output to a finite candidate set
$(\candset, \probdist)$ via the LLM's self-reported preference
weights. Single call, $K{=}4$ default. Non-compliant return
(weights missing / $|\candset|{<}2$) falls back to a fresh
single-answer call that bypasses sampling and does not log an audit.}
\label{fig:prompt-template}
\end{figure}


% ============================================================
\subsection{Self-Reported Candidate Sampling}
\label{sec:method-sampler}
% ============================================================
\todo{Parse JSON into $(\text{decisions}, \text{weights})$, normalize
weights to $\probdist$, feed to \agentmark's binning sampler with
keyed context $\ctx \| \nonceT$ where $\nonceT = \text{PRF}(K,
\ctx)$. The sampler returns the keyed pick $c^*$ and embedded bits.
Algorithm \ref{alg:sample} sketches the per-call flow.}

% ============================================================
% Algorithm: per-call keyed sample
% ============================================================
\begin{algorithm}[t]
\caption{Watermarked LLM call (per-decision)}
\label{alg:sample}
\begin{algorithmic}[1]
\Require Original SDK prompt $P$; backend LLM client $L$; PRF key $K$;
target candidate count $K_{\text{cands}}{=}4$; context payload $\ctx$
\Ensure SDK-schema-conformant string $c^*$; audit record $\reveal_t$
\Statex
\State $P' \gets P \,\|\, \textsc{AgentMarkInstruction}(K_{\text{cands}})$
\State $\text{raw} \gets L.\text{complete}(P')$ \Comment{1 call, $T{=}0.7$}
\State $(\text{decisions}, \text{weights}) \gets \textsc{ParseJSON}(\text{raw})$
\If{$|\text{decisions}| < 2$ \textbf{or} $\text{weights missing}$}
  \State \Return $L.\text{complete}(P)$ \Comment{passthrough, no audit}
\EndIf
\State $\probdist \gets \textsc{Normalize}(\text{weights})$
\State $\nonceT \gets \text{PRF}(K, \ctx)$
\State $c^* \gets \textsc{BinningSample}(\text{decisions}, \probdist, \nonceT)$
\State $\commit \gets H(\ctx \| H(\candset) \| H(\probdist) \| \text{idx}(c^*) \| \text{bits} \| \nonceT)$
\State $\textsc{MerkleLog}.\text{append}(\commit)$
\State $\reveal_t \gets \{\candset, \probdist, c^*, \nonceT, \ctx\}$
\State \Return $(c^*, \reveal_t)$
\end{algorithmic}
\end{algorithm}


\paragraph{Properties.} \todo{State Lemma 1 (strict
distribution preservation), Lemma 2 (composition across the cascade),
Lemma 3 (backend-invariant marginal). Full proofs in
Appendix~\ref{app:proofs}.}

\begin{lemma}[Strict distribution preservation]
\label{lem:dist-preserve}
\todo{Given $\probdist$, the keyed pick $c^*$ under PRF key $K$ has
marginal distribution equal to $\probdist$: $\Pr[c^* = c_i \mid K,
\ctx]_{\mathrm{marg}} = \probdist(c_i)$.}
\end{lemma}

\begin{lemma}[Cascade composition]
\label{lem:cascade}
\todo{$K$ LLM calls within a memory event yield independent keyed
picks. Capacity is additive in $K$, marginal bias does not
accumulate.}
\end{lemma}

\begin{lemma}[Backend-invariant marginal]
\label{lem:backend-invariant}
\todo{Lemma \ref{lem:dist-preserve} depends only on $\probdist$ and
$K$, not the backend.}
\end{lemma}

% ============================================================
\subsection{Cryptographic Audit Trace}
\label{sec:method-audit}
% ============================================================
\todo{Three-layer design: per-decision commitment, per-session
Merkle log, signed root anchor written into the snapshot itself.
$\commit = H(\ctx \| H(\candset) \| H(\probdist) \| \text{selected}
\| \text{bits} \| \nonceT)$. $\header_T = (\text{agent\_id},
\text{user\_id}, \text{session\_id}, T, \text{root}_T,
\text{sig}_K(\text{root}_T))$.}

\paragraph{Per-leaf inclusion proofs.}
\todo{Each \texttt{AuditRecord} stores its Merkle inclusion path at
seal time so the verifier can check each leaf independently against
$\text{root}_T$ — no rebuilt-root match required. This gives smooth
bit-recovery under structural attacks (pruning / dedup / poisoning),
in contrast to rebuilt-root verification which binary-collapses when
the leaf set changes (see Table~\ref{tab:robustness}).}

\paragraph{Verification regimes.}
\todo{Define R1 (full external log), R2 (partial external log), R3
(snapshot only). Forward to \S\ref{sec:rq3} for the empirical
hierarchy and to Table~\ref{tab:verification} for the regime
matrix.}
